\section{Positioning: what may be claimed}\label{sec:positioning}

\subsection{Retracted claims}\label{sec:retracted}

The following statements appeared in earlier drafts of this project and are
withdrawn. They are recorded here rather than silently deleted, because a
retraction that is not written down tends to reappear in the next draft.

\begin{enumerate}[label=R\arabic*.,leftmargin=*]
\item \emph{``First complete closed-form realization of the two revised-simplex
      basis systems for a forest-structured graph LP.''} Withdrawn: too close to
      \citet{Yang2026}, who develops exactly that on a budget-constrained
      total-variation program.
\item \emph{``First closed-form forest simplex.''} Withdrawn: too broad, and
      contradicted both by \citet{Yang2026} and by the classical
      network-simplex-with-side-constraint literature
      \citep{AhujaEtAl1993,Schrijver2003}.
\item \emph{Novelty of the forest structure of basic feasible solutions.}
      Withdrawn: present in \citet{Wuille2025} and formalised in
      \citet{MollakhaniEtAl2026} for the same normalised program.
\item \emph{Novelty of the merge/split reading of pivots.} Withdrawn: same
      sources, and \citet{Yang2026} independently.
\item \emph{Novelty of the forest state and of the aggregate sign test, even
      relative to the 2025--2026 work.} Withdrawn, and further back than the
      previous two items: \citet{LerchsGrossmann1965} maintain a forest of
      branches on the closure graph, carry the mass of each branch, classify an
      arc by the sign of the mass it supports, and pivot by adding one arc and
      removing the branch's root arc (\cref{sec:lg}). With $b_i = p_i - \rho
      w_i$ their strong/weak test is the sign of the \FS{} reduced cost. Nothing
      about maintaining a forest and testing aggregates on it can be claimed.
\item \emph{Novelty of the maximum-ratio closure formulation, of its solution by
      max-flow, and of the canonical decomposition.} Withdrawn: classical
      \citep{Picard1976,PicardQueyranne1982,Sidney1975,Lawler1978}.
\item \emph{Novelty of the ``simplex on an incidence matrix plus one side row''
      paradigm.} Withdrawn, and datable. \citet{ChenSaigal1977} prove that a
      basis of a capacitated flow problem with $k$ additional linear rows is a
      spanning tree plus $k$ columns, and drive pricing, pivoting and inversion
      from a working basis of size $k$; \citet{Caliskan2011} specialises the
      case $k=1$ to maximum flow with a budget row, where the basis is a
      spanning tree plus one arc closing a fundamental cycle of nonzero net
      cost; \citet{HolzhauserEtAl2017} do the same for minimum cost flow. The
      paradigm is due to 1977, not to the recent literature.
\item \emph{Novelty of solving the ratio version by walking a forest state.}
      Withdrawn: \citet{Hochbaum2001} gives a parametric Lerchs--Grossmann with
      the complexity of a single run, and \citet{PicardSmith2004} compute
      open-pit contours maximising a benefit--cost ratio by parametric flow.
\item \emph{Novelty of the equality normalisation \(\sum_i w_i y_i = 1\).}
      Withdrawn: it is the homogeneous case of the reduction of
      \citet{CharnesCooper1962}, whose Theorem~2 already writes the
      normalisation as an equality with right-hand side $1$
      (\cref{sec:open}, Q5).
\item \emph{Any claim of superior practical performance.} Withdrawn on the
      project's own evidence: on the frozen test bed the parametric and
      Dinkelbach max-flow backends dominate \FS{} on both solved count and time
      (\cref{sec:open}).
\end{enumerate}

\subsection{The cell nobody occupies}\label{sec:gap}

The claim reduces to one empty cell in a table with two axes: what the method
maintains, and which program it solves.

\begin{center}
\begin{tabular}{@{}lll@{}}
\toprule
& maximum closure, weights fixed & maximum-ratio closure \\
\midrule
flow / cut state
  & \citet{Picard1976}, \citet{Hochbaum2008}
  & \citet{GalloEtAl1989}, \citet{Hochbaum2024} \\
forest state, pivots, no LP basis
  & \citet{LerchsGrossmann1965}
  & \citet{Hochbaum2001} (parametric LG), \citet{MollakhaniEtAl2026} \\
exact simplex, LP basis
  & ---
  & \textbf{this work} \\
\addlinespace
\multicolumn{3}{@{}l@{}}{\emph{exact simplex on a sibling program:}
  \citet{Yang2026} (budget-constrained total variation),
  \citet{HolzhauserEtAl2017} (budget-constrained min cost flow),} \\
\multicolumn{3}{@{}l@{}}{\quad
  \citet{ChenSaigal1977} ($k$ side rows on a capacitated flow),
  \citet{Caliskan2011} (budget-constrained max flow)} \\
\bottomrule
\end{tabular}
\end{center}

Read across the bottom row: the exact simplex has been developed on programs
in the same structural family as \eqref{eq:class} --- a difference system or a
network with one budget row --- but not on \eqref{eq:class} itself. Read up the
right column: the ratio program has been solved by flow states and by forest
states that pivot, but not by a simplex maintaining an LP basis with reduced
costs, a ratio test and dual multipliers. \citet{MollakhaniEtAl2026} say so
themselves: their algorithm is, in the abstract's own words, motivated by
structural properties of basic feasible solutions of the simplex method, and
what it then maintains and refines is a partition of the vertices into
dependency-respecting subsets ordered by decreasing aggregate ratio --- chunks
merged and split directly on the graph, not a basis.

\paragraph{The most serious candidate for the empty cell, and why it does not
fill it.} \citet{MargotEtAl2003} is the one paper found in this review that
promised to occupy the bottom-right cell: an in-depth theoretical, algorithmic
and computational study of an LP relaxation of a precedence-constrained problem,
whose algorithmic core is a parametric extension of Sidney's decomposition and
whose finest decomposition is obtained, in their words, by essentially solving a
parametric minimum-cut problem. Read in full, it belongs in the first row and
not in the third: the state it maintains is a sequence of nested minimum cuts,
and the only simplex in the paper is the one inside the general-purpose LP
solver used to produce the computational results. It carries no basis
characterisation, no reduced costs and no ratio test for the relaxation. It
strengthens R6 and R8 --- the decomposition and its parametric reading are
classical --- and leaves the cell empty.

\paragraph{A warning that belongs next to the table.} An empty cell is not by
itself a contribution. It can be empty because the route is worse, and the
evidence available here points that way: on this project's own test bed the
flow-based entries of the first row and the parametric entry of the second
dominate the bottom row on both solved count and time (\cref{sec:open}). The
honest reading is that this work fills the cell and reports what is inside it,
including the part that is unfavourable.

\subsection{What is left}

Four candidates survive the review. None of them is a paradigm; all of them are
specialisations, and each carries a condition that has not yet been discharged.

\begin{enumerate}[label=C\arabic*.,leftmargin=*]
\item \textbf{Specialisation of the ordinary simplex, primal and dual, to the
      normalised maximum-ratio closure LP, and to the canonical decomposition
      it drives.} That a simplex method is exact is not a contribution and is
      not offered as one; what is at stake is what the method \emph{becomes}
      on this program. \citet{Wuille2025} and
      \citet{MollakhaniEtAl2026} expose the forest structure of the basic
      feasible solutions of this program and then leave the simplex, optimising
      the linearization directly; \citet{Yang2026} develops the exact simplex,
      but on a different program; \citet{LerchsGrossmann1965} pivot on a forest
      of the closure graph, but outside the simplex and at fixed weights.
      Carrying the ordinary simplex through on this LP --- basis state, prices,
      pivot directions, ratio test, multipliers, certificates, in a primal and
      in a dual variant --- and then driving the whole canonical sequence of
      blocks with it, not merely the first, is the gap between all three.
      \emph{Condition:} it must be presented as a specialisation, not as a new
      method; the debt to all three lines must be stated in the abstract, not
      only in related work; and the difference from a normalized tree must be
      made precise rather than asserted, because to a mining audience the two
      objects look identical.

\item \textbf{The closure-specific closed forms are markedly simpler than the
      general forest-simplex algebra.} For the normalised program the reduced
      costs collapse to
      \[
        \bar c_{y_T} = p(T) - \rho\,w(T),
        \qquad
        \bar c_{s_e} = \sigma_e\bigl[p(D_e) - \rho\,w(D_e)\bigr],
      \]
      with no edge objective coefficients, no box-bound contributions, no
      orientation choice between positive and negative edge variables and no
      generic working-basis algebra. \emph{Condition:} the comparison against
      \citet{Yang2026} must be made explicitly and term by term; ``simpler''
      is not a claim until the two sets of formulas are put side by side.

\item \textbf{Full pricing in $O(n)$ arithmetic, independent of the number of
      precedence arcs.} There are exactly $n-1$ nonbasic candidates; all forest
      aggregates are obtained in one postorder; each candidate costs $O(1)$; no
      non-forest arc is ever priced. \emph{Condition:} this is the cleanest of
      the four, and also the one most exposed. If \citet{Yang2026} admits an
      equivalent $O(n)$ pricing in his model, the claim is not that $O(n)$
      pricing is possible but that it is available here without the structural
      overhead his model carries. That check has not been done.

\item \textbf{The separation between pricing and feasibility in the exact ratio
      test.} Pricing ranges over the $n-1$ nonbasic forest and root variables;
      feasibility ranges over all potentially decreasing basic slacks. The
      distinction is clean and it is what makes the exact simplex correct rather
      than merely plausible. \emph{Condition:} on its own this is a component of
      C1, not an independent contribution, and should be presented as such.
\end{enumerate}

\subsection{The formulation that survives}

\begin{quote}
To the best of our knowledge we give the first explicit simplex realisation,
primal and dual, specialised to the normalised maximum-ratio closure LP, and we
use it to produce the entire canonical sequence of blocks rather than the
maximum-ratio block alone. Building on the forest structure of its basic
feasible solutions --- which is closely related to recent spanning-forest
linearization work --- we derive problem-specific closed forms for all nonbasic
prices and pivot directions. In particular, complete pricing requires $O(n)$
arithmetic independently of the number of precedence arcs, while a full ratio
test accounts for all potentially blocking basic slacks. The canonical
decomposition itself is classical and is not claimed.
\end{quote}

\noindent followed immediately, and in the same place, by:

\begin{quote}
We do not claim novelty for the forest structure itself, for the merge/split
basis effects, for the maximum-ratio closure formulation, for the canonical
decomposition, or for the general paradigm of simplex methods on incidence
matrices with side constraints; these have direct precedents in
\citet{Wuille2025} and \citet{MollakhaniEtAl2026}, in \citet{Yang2026}, and in
the classical network-simplex literature.
\end{quote}

\subsection{The name}\label{sec:name}

The name follows the positioning, and the earlier one contradicted it. ``Forest
Simplex'' named the forest --- the one element \cref{sec:retracted} withdraws
--- and it collided three ways: with the \emph{spanning forest} of
\citet{MollakhaniEtAl2026} and \citet{Wuille2025}, with the \emph{rooted
spanning forests} of \citet{Yang2026}, and, one letter and no syllables away,
with the Forrest--Tomlin update of the basis factors, which is what an audience
of linear programmers hears first.

The method is therefore called the \FS{}, lower case, on the model of
\emph{network simplex}: a name taken from the problem the method is specialised
to, not from the structure it maintains. The problem is
\(\max_C p(C)/w(C)\) over the closures of a digraph, and the linear program is
\eqref{eq:class}; the capacity row of the application is absent from both,
which is why it is absent from the name as well. The precedence-constrained knapsack is where
the program comes from and where the method is measured, and it is reported as
such among the applications; it is kept out of the title as well as out of the
name, because the method delivers the whole canonical sequence and not only the
first block, and a knapsack in the title would promise that the capacity row is
handled by the simplex, which it is not.
The forest keeps its own name, as the structure of the basis, with the
attribution \cref{sec:lg} gives it.

\subsection{The question the work answers}

The question is not whether maximum-ratio closure can be solved with forests ---
it can, and that is known. It is:

\begin{quote}
what does the exact ordinary primal simplex look like on the normalised
maximum-ratio closure LP, once the known forest structure of its basic feasible
solutions is exploited completely?
\end{quote}

The answer is the algorithm: basis state as forest plus root, exact aggregate
prices, $O(n)$ pricing, exact forest motion, full ratio test, dual multipliers
and certificates, and an implementation measured against the alternatives. That
last item is where the work currently stands weakest, and \cref{sec:open} says
why.
